% Appendix A

\chapter{Data Appendix}\label{app:data}
\begin{lstlisting}[caption={Example file}, label={lst:example_file}]
	_exptl_crystal_grow.crystal_id      1
	_exptl_crystal_grow.method          'VAPOR DIFFUSION, SITTING DROP'
	_exptl_crystal_grow.apparatus       None
	_exptl_crystal_grow.atmosphere      None
	_exptl_crystal_grow.pH              9.5
	_exptl_crystal_grow.temp            298.0
	_exptl_crystal_grow.pdbx_details    '3.2M (NH4)2SO4, 0.1M Glycine, ...'
	_exptl_crystal_grow.time            None
\end{lstlisting}

\hspace{0.5cm}
\begin{lstlisting}[caption={Feature example}, label={lst:feature_example}]
	loop_
	_atom_site.group_PDB
	_atom_site.id
	_atom_site.type_symbol
	_atom_site.label_atom_id
	_atom_site.label_alt_id
	_atom_site.label_comp_id
	_atom_site.label_asym_id
	_atom_site.label_entity_id
	_atom_site.label_seq_id
	_atom_site.pdbx_PDB_ins_code
	_atom_site.Cartn_x
	_atom_site.Cartn_y
	_atom_site.Cartn_z 
	ATOM   1    N  N   . VAL A 1 1   ? 6.204   16.869  4.854   1.00 49.05 ...
	ATOM   2    C  CA  . VAL A 1 1   ? 6.913   17.759  4.607   1.00 43.14 ...
	ATOM   3    C  C   . VAL A 1 1   ? 8.504   17.378  4.797   1.00 24.80 ...
	ATOM   4    O  O   . VAL A 1 1   ? 8.805   17.011  5.943   1.00 37.68 ...
\end{lstlisting}

\begin{figure}[h!]
	\centering		
	\begin{subfigure}{0.489\linewidth}
		\centering
		\includegraphics[width=\linewidth]{../msc-thesis-code/data/dataExploration/plots/phGivenVsParsedpH}
		\caption[Parsed pH vs. entered pH]{ }
		\label{fig:enteredVsParsedPh}
	\end{subfigure}
	\hspace{0.1cm}
	\begin{subfigure}[b]{0.489\linewidth}
		\centering
		\includegraphics[width=\linewidth]{../msc-thesis-code/data/dataExploration/plots/tempGivenVsParsedTemp}
		\caption[Parsed temperature vs. entered temperature]{ }
		\label{fig:enteredVsParsedTemp}
	\end{subfigure}
	\caption[Parsing Quality for numerical values]{\autoref{fig:enteredVsParsedPh} shows the pH values parsed from free text with the single pH value recorded in the numerical field. Multiple pH values reported in free text for different solutions leads to mismatches. In contrast, the temperature field for crystallization produced almost no mismatches since only one value was given in free text as seen in \autoref{fig:enteredVsParsedPh}.}
	\label{fig:enteredVsParsed}
\end{figure}


\begin{figure}[h!]
	\centering		
	\begin{subfigure}[b]{0.495\linewidth}
		\centering
		\includegraphics[width=\linewidth]{../msc-thesis-code/data/dataExploration/plots/gravyHist}
		\caption[Kyte-Doolittle hydropathy distribution]{ }
		\label{fig:gravyHist}
	\end{subfigure}
	\begin{subfigure}[b]{0.495\linewidth}
		\centering
		\includegraphics[width=\linewidth]{../msc-thesis-code/data/dataExploration/plots/aromHist}
		\caption[Aromaticity distribution]{ }
		\label{fig:aromaticity}
	\end{subfigure}
	\begin{subfigure}[b]{0.495\linewidth}
		\centering
		\includegraphics[width=\linewidth]{../msc-thesis-code/data/dataExploration/plots/pIHist}
		\caption[Isoelectric point distribution]{ }
		\label{fig:pIHist}
	\end{subfigure}
	\begin{subfigure}[b]{0.495\linewidth}
		\centering
		\includegraphics[width=\linewidth]{../msc-thesis-code/data/dataExploration/plots/mwHist}
		\caption[Molecular weight distribution]{ }
		\label{fig:mwHist}
	\end{subfigure}
	\caption[Distributions of sequence-derived features]{Distributions of sequence-derived features across the dataset. 
		Kyte--Doolittle hydropathy (\autoref{fig:gravyHist}) scores are roughly normally distributed with a slightly negative mean, while aromaticity (\autoref{fig:aromaticity}) values form a mildly positive normal distribution. 
		The isoelectric point (\autoref{fig:pIHist}) shows a bimodal pattern with peaks near pH~6 and pH~9. 
		Molecular weights (\autoref{fig:mwHist}) decrease sharply beyond 50{,}000~Da, mirroring trends in sequence length and atom count seen in \autoref{fig:sequenceAndAtomCount}.}
	
	\label{fig:sequenceInformation}
\end{figure}

\begin{figure}[h!]
	\centering		
	\begin{subfigure}[b]{0.495\linewidth}
		\centering
		\includegraphics[width=\linewidth]{../msc-thesis-code/data/dataExploration/plots/temperatureDistribution}
		\caption[Temperature Distribution]{ }
		\label{fig:temperatureDistrbution}
	\end{subfigure}
	\begin{subfigure}[b]{0.495\linewidth}
		\centering
		\includegraphics[width=\linewidth]{../msc-thesis-code/data/dataExploration/plots/phDistribution}
		\caption[pH Distribution]{ }
		\label{fig:pHDistribution}
	\end{subfigure}
	\caption[Temperature and pH distributions]{Distribution of numerical fields in the condition category. Only a few distinct temperatures (\ref{fig:temperatureDistrbution}) are reported, with 35\% of samples specifying room temperature and most remaining values clustered around it; a secondary peak appears at 4\,$^\circ$C. In contrast, pH values follow an approximately normal distribution centered near pH\roughly7 (\ref{fig:pHDistribution}).}
	\label{fig:temperatureAndPh}
\end{figure}


\begin{figure}[h!]
	\centering		
	\begin{subfigure}[t]{0.489\linewidth}
		\centering
		\includegraphics[width=\linewidth]{../msc-thesis-code/data/dataExploration/plots/classifiedConcentration}
		\caption[Unit specification]{ }
		\label{fig:classifiedConcentration}
	\end{subfigure}
	\hspace{0.1cm}
	\begin{subfigure}[t]{0.489\linewidth}
		\centering
		\includegraphics[width=\linewidth]{../msc-thesis-code/data/dataExploration/plots/weight_per_volume_specification_ratio}
		\caption[Ratio of unit specifications against Molecular weight]{ }
		\label{fig:weight_per_volume_specification_ratio}
	\end{subfigure}
	\caption[\gls{peg} Unit analysis]{\autoref{fig:classifiedConcentration} shows the distribution of named units. For the majority of the parsed PEG concentration no unit was specified. Concentration can be given both in weight per volume and volume per volume. This naming is related to the molecular weight as seen in \autoref{fig:weight_per_volume_specification_ratio}, which shows the relation between molecular weight and how often the unit was specified to be $w/v$ in relation to all named unit appearances.}
	\label{fig:unitSpecification}
\end{figure}

\begin{figure}[h!]
	\centering		
	\begin{subfigure}[t]{0.49\linewidth}
		\centering
		\includegraphics[width=\linewidth]{../msc-thesis-code/data/dataExploration/plots/aminoAcidOccurence}
		\caption[Residue distribution]{ }
		\label{fig:aminoAcidOccurence}
	\end{subfigure}
	\hspace{0.1cm}
	\begin{subfigure}[t]{0.49\linewidth}
		\centering
		\includegraphics[width=\linewidth]{../msc-thesis-code/data/dataExploration/plots/observedVsPredictedAminoAcidCount}
		\caption[Observed vs. Predicted Occurrences of Residues]{ }
		\label{fig:observedVsPredictedAminoAcidCount}
	\end{subfigure}
	\caption[Residue analysis in protein sequences]{
		\autoref{fig:aminoAcidOccurence} shows the distribution of amino acids across all sequences. Residue frequencies vary substantially—Leucine is over four times more common than Tryptophan. While \autoref{fig:observedVsPredictedAminoAcidCount} compares observed residue frequencies with those predicted from codon usage.}
	
	\label{fig:aminoAcidCounts}
\end{figure}


\begin{figure}[h]
	\centering
	\includegraphics[width=1\linewidth]{../msc-thesis-code/data/dataExploration/plots/employmentMethodDistribution}
	\caption[Most common crystallization methods]{A visualization of the most common methods employed to crystallize a protein. There is virtually no variation in method with more than 95\% of entries being crystallized using a variation of vapor diffusion.}
	
	\label{fig:employmentMethodDistribution}
\end{figure}

\begin{figure}[H]
	\centering
	\includegraphics[width=0.6\linewidth]{../msc-thesis-code/data/dataExploration/plots/correlation_ph_temp}
	\caption[Temperature / pH correlation with surface features]{Correlation matrix displaying correlation factors for the temperature and pH. Label values do not exhibit strong correlations directly with any of the features. However, there are some correlation between surface features as well.}
	\label{fig:correlationphtemp}
\end{figure}

\begin{figure}[H]
	\centering
	\includegraphics[width=0.6\linewidth]{../msc-thesis-code/data/dataExploration/plots/db_code}
	\caption[Number of unique counts per protein]{The figure shows the unique number of entries per protein. Some proteins have been crystallized extremely often.}
	\label{fig:most_common_proteins}
\end{figure}


\begin{figure}[h!]
	\centering		
	\begin{subfigure}[t]{0.49\linewidth}
		\centering
		\includegraphics[width=\linewidth]{../msc-thesis-code/data/dataExploration/plots/net_charge_normalizedVsChem}
		\caption[Distributions normalized surface net charge]{ }
		\label{fig:net_charge_normalizedVsChem}
	\end{subfigure}
	\hspace{0.1cm}
	\begin{subfigure}[t]{0.49\linewidth}
		\centering
		\includegraphics[width=\linewidth]{../msc-thesis-code/data/dataExploration/plots/surface_entropy_fractionVsChem}
		\caption[Distributions surface entropy]{ }
		\label{fig:hydrophobic_surface_fractionVsChem}
	\end{subfigure}
	\begin{subfigure}[t]{0.49\linewidth}
		\centering
		\includegraphics[width=\linewidth]{../msc-thesis-code/data/dataExploration/plots/sasa_skewnessVsChem}
		\caption[Distributions SASA skewness]{ }
		\label{fig:sasa_skewnessVsChem}
	\end{subfigure}
	\hspace{0.1cm}
	\begin{subfigure}[t]{0.49\linewidth}
		\centering
		\includegraphics[width=\linewidth]{../msc-thesis-code/data/dataExploration/plots/sidechain_fractionVsChem}
		\caption[Distributions side chain fraction]{ }
		\label{fig:sidechain_fractionVsChem}
	\end{subfigure}
	\begin{subfigure}[t]{0.49\linewidth}
		\centering
		\includegraphics[width=\linewidth]{../msc-thesis-code/data/dataExploration/plots/roughnessVsChem}
		\caption[Distributions normalized roughness]{ }
		\label{fig:roughnessVsChem}
	\end{subfigure}
	\hspace{0.1cm}
	\begin{subfigure}[t]{0.49\linewidth}
		\centering
		\includegraphics[width=\linewidth]{../msc-thesis-code/data/dataExploration/plots/hydro_polar_contrastVsChem}
		\caption[Distributions hydro polar contrast]{ }
		\label{fig:hydro_polar_contrastVsChem}
	\end{subfigure}
	\caption[Surface feature distributions]{
		Distributions of different chemicals for different surface features}
	\label{fig:surfaceFeatureChemicalDistributions}
\end{figure}

\begin{figure}[h!]
	\centering		
	\begin{subfigure}[t]{0.49\linewidth}
		\centering
		\includegraphics[width=\linewidth]{../msc-thesis-code/data/dataExploration/plots/net_charge_normalized_ks_distance}
		\caption[Distributions normalized surface net charge]{ }
		\label{fig:net_charge_normalized_ks_distance}
	\end{subfigure}
	\hspace{0.1cm}
	\begin{subfigure}[t]{0.49\linewidth}
		\centering
		\includegraphics[width=\linewidth]{../msc-thesis-code/data/dataExploration/plots/surface_entropy_fraction_ks_distance}
		\caption[Distributions surface entropy]{ }
		\label{fig:hydrophobic_surface_fraction_ks_distance}
	\end{subfigure}
	\begin{subfigure}[t]{0.49\linewidth}
		\centering
		\includegraphics[width=\linewidth]{../msc-thesis-code/data/dataExploration/plots/sasa_skewness_ks_distance}
		\caption[Distributions SASA skewness]{ }
		\label{fig:sasa_skewness_ks_distance}
	\end{subfigure}
	\hspace{0.1cm}
	\begin{subfigure}[t]{0.49\linewidth}
		\centering
		\includegraphics[width=\linewidth]{../msc-thesis-code/data/dataExploration/plots/sidechain_fraction_ks_distance}
		\caption[Distributions side chain fraction]{ }
		\label{fig:sidechain_fraction_ks_distance}
	\end{subfigure}
	\begin{subfigure}[t]{0.49\linewidth}
		\centering
		\includegraphics[width=\linewidth]{../msc-thesis-code/data/dataExploration/plots/roughness_ks_distance}
		\caption[Distributions normalized roughness]{ }
		\label{fig:roughness_ks_distance}
	\end{subfigure}
	\hspace{0.1cm}
	\begin{subfigure}[t]{0.49\linewidth}
		\centering
		\includegraphics[width=\linewidth]{../msc-thesis-code/data/dataExploration/plots/hydro_polar_contrast_ks_distance}
		\caption[Distributions hydro polar contrast]{ }
		\label{fig:hydro_polar_contrast_ks_distance}
	\end{subfigure}
	\caption[Surface feature distributions]{
		Distributions of different chemicals for different surface features}
	\label{fig:surfaceFeatureKsDistance}
\end{figure}
